\documentclass[11pt]{article}
\usepackage[margin=1in]{geometry}
\usepackage{amsmath,amssymb}
\usepackage{hyperref}
\usepackage{booktabs}

\title{Rapidity shape of the Pb--Pb cross section:\\
why our peak is at more negative $y$ than in arXiv:2606.05469}
\author{}
\date{\today}

\begin{document}
\maketitle

\section{Summary}

\texttt{plotting\_scripts/PbPb\_bins.py} plots our exclusive + diffractive
$D^0$ cross section in Pb--Pb (An0n) in the bins of Fig.~6 of Guzey,
Innocenti, Sta\'sto and Strikman (arXiv:2606.05469). At high $p_T$ the two
shapes in rapidity are different:
\begin{itemize}
  \item In the paper the bin $-1<y<0$ is clearly the highest.
  \item In our result the bin $-2<y<-1$ is as high as, or higher than, the bin
    $-1<y<0$. Our peak is about half a unit of rapidity further towards
    negative $y$.
\end{itemize}
The most likely reason is the region $x_\mathbb{P}>0.01$. Negative rapidity
needs large $x_\mathbb{P}$. The paper suppresses large $x_\mathbb{P}$ with the
nuclear form factor; our calculation has no suppression there. This is tested
only at $p_{D^0\perp}=2$~GeV (Sec.~\ref{sec:test}), so it is a strong hint and
not a proof.

\section{The difference in numbers}

Bin averages of $d\sigma/dy\,dp_{D^0\perp}$ in mb/GeV, BCFY, An0n. The values
of the paper are read off the ``high shadowing'' band of its Fig.~6 and are
approximate.
\begin{center}
\begin{tabular}{lcccc}
\toprule
 & \multicolumn{2}{c}{ours (central value)} & \multicolumn{2}{c}{arXiv:2606.05469} \\
$p_{D^0\perp}$ bin & $-2<y<-1$ & $-1<y<0$ & $-2<y<-1$ & $-1<y<0$ \\
\midrule
2--3 GeV   & $0.150$ & $0.149$ & $\sim0.68$ & $\sim0.73$ \\
6--7 GeV   & $6.6\times10^{-4}$ & $6.2\times10^{-4}$ & $\sim6.4\times10^{-3}$ & $\sim8.4\times10^{-3}$ \\
10--11 GeV & $5.24\times10^{-5}$ & $5.20\times10^{-5}$ & $\sim3.3\times10^{-4}$ & $\sim4.8\times10^{-4}$ \\
\bottomrule
\end{tabular}
\end{center}
At 2--3 GeV the two shapes are similar. The difference grows with
$p_{D^0\perp}$.

\section{The two parts of our result}

Our cross section at $p_{D^0\perp}=10.5$~GeV (BCFY, An0n), in units of
$10^{-5}$~mb/GeV:
\begin{center}
\begin{tabular}{lccccccc}
\toprule
$y$ & $-2.5$ & $-2$ & $-1.5$ & $-1$ & $-0.5$ & $0$ & $0.5$ \\
\midrule
diffractive & 1.8 & 3.6 & 5.0 & 5.5 & 5.0 & 3.8 & 2.4 \\
exclusive   & 0.42 & 0.42 & 0.37 & 0.31 & 0.27 & 0.24 & 0.19 \\
\bottomrule
\end{tabular}
\end{center}
\begin{itemize}
  \item The diffractive part has its maximum at $y=-1$ and is symmetric around
    it. The bins $(-2,-1)$ and $(-1,0)$ are then equal.
  \item The exclusive part keeps growing towards negative $y$. It makes the
    sum slightly larger in $(-2,-1)$.
\end{itemize}
In the paper the maximum is at about $y=-0.5$.

\section{Negative rapidity needs large $x_\mathbb{P}$}

In the code (\texttt{src/integrands.cpp}) the charm quark has
$p^+=m_T e^{y}/\sqrt2$, the photon has $q^+=p^+/z$, and
\begin{equation}
  W^2=\sqrt2\,q^+\sqrt{s_{NN}}, \qquad
  M_{q\bar q}^2=\frac{m_T^2}{z(1-z)}, \qquad
  x_\mathbb{P}=\frac{M_{q\bar q}^2}{\beta\,W^2}
  =\frac{m_T\,e^{-y}}{\beta\,(1-z)\sqrt{s_{NN}}} ,
\end{equation}
with $\beta=1$ for the exclusive part (in the code $\beta$ is called
\texttt{x}). Here $m_T$ is the transverse mass of the charm quark,
$m_T^2=p_{D^0\perp}^2/z_h^2+m_c^2$. So $x_\mathbb{P}$ grows as $e^{-y}$ and
linearly with $m_T$.

For $z=1/2$ and $\beta=1$ this is $x_\mathbb{P}=2m_Te^{-y}/\sqrt{s_{NN}}$. As
an example take $p_{D^0\perp}=10.5$~GeV and $z_h\approx0.7$, so
$m_T\approx15$~GeV, with $\sqrt{s_{NN}}=5.36$~TeV:
\begin{center}
\begin{tabular}{lcccc}
\toprule
$y$ & $0$ & $-1$ & $-2$ & $-3$ \\
\midrule
$x_\mathbb{P}$ & $0.006$ & $0.015$ & $0.04$ & $0.11$ \\
\bottomrule
\end{tabular}
\end{center}
For the diffractive part $\beta<1$ and $x_\mathbb{P}$ is larger than this.
At high $p_T$ and $y<-1$ all the cross section comes from
$x_\mathbb{P}>0.01$.

\section{How large $x_\mathbb{P}$ is treated}

\paragraph{Our code.} The dipole amplitude is known for $x<0.01$. For larger
$x_\mathbb{P}$ the code uses the amplitude at $x=0.01$
(\texttt{min(x\_po, 0.01)}). There is no nuclear form factor for the
longitudinal momentum transfer. The only limit is the cut
$x_\mathbb{P}<0.1$. So nothing reduces the cross section between
$x_\mathbb{P}=0.01$ and $0.1$.

\paragraph{arXiv:2606.05469.} The nuclear diffractive PDFs (their Eq.~8)
have the factor $e^{i z' x_\mathbb{P} m_N}$ inside the integral over the
nuclear density. This is the nuclear form factor at
$|t_{\min}|\simeq x_\mathbb{P}^2m_N^2$: the nucleus stays intact only if the
longitudinal momentum transfer $x_\mathbb{P}m_N$ is small compared with
$1/R_A$. The paper also says that the cross section decreases at large
negative rapidity because larger momentum fractions are probed there.

\paragraph{Size of the form factor.} For lead, with
\begin{equation}
  F_A(q)=\frac{3\left[\sin(qR_A)-qR_A\cos(qR_A)\right]}{(qR_A)^3}\,
  \frac{1}{1+a^2q^2}, \qquad q=x_\mathbb{P}m_N,
  \label{eq:FA}
\end{equation}
$R_A=6.62$~fm and $a=0.7$~fm:
\begin{center}
\begin{tabular}{lcccccc}
\toprule
$x_\mathbb{P}$ & 0.01 & 0.02 & 0.03 & 0.05 & 0.07 & 0.1 \\
\midrule
$F_A^2$ & 0.98 & 0.92 & 0.82 & 0.57 & 0.32 & 0.07 \\
\bottomrule
\end{tabular}
\end{center}

\section{Test with our $x_\mathbb{P}$-differential results}
\label{sec:test}

\texttt{run\_xpom.sh} gives $d\sigma/d\ln x_\mathbb{P}$ at
$p_{D^0\perp}=2$~GeV (BCFY, An0n; \texttt{output/An0n/xpom/}). The table shows
the integral over $x_\mathbb{P}$, the part of it that comes from
$x_\mathbb{P}>0.01$, and the integral after multiplying the integrand by
$F_A^2$ of Eq.~\eqref{eq:FA}. Cross sections are in mb/GeV.
\begin{center}
\begin{tabular}{lccccc}
\toprule
 & $y$ & total & from $x_\mathbb{P}>0.01$ & with $F_A^2$ & ratio \\
\midrule
diffractive & $-2$ & 0.0523  & 99\% & 0.0215  & 0.41 \\
            & $-1$ & 0.0508  & 97\% & 0.0237  & 0.47 \\
            & $0$  & 0.0378  & 94\% & 0.0188  & 0.50 \\
            & $1$  & 0.0211  & 92\% & 0.0109  & 0.52 \\
            & $2$  & 0.00754 & 91\% & 0.0040  & 0.53 \\
            & $3$  & 0.00126 & 90\% & 0.00067 & 0.54 \\
\midrule
exclusive   & $-2$ & 0.269  & 40\% & 0.244  & 0.91 \\
            & $-1$ & 0.281  & 13\% & 0.273  & 0.97 \\
            & $0$  & 0.272  & 4\%  & 0.270  & 0.99 \\
            & $1$  & 0.210  & 1\%  & 0.210  & 1.00 \\
            & $2$  & 0.106  & 0.4\% & 0.106 & 1.00 \\
            & $3$  & 0.0242 & 0.1\% & 0.0242 & 1.00 \\
\bottomrule
\end{tabular}
\end{center}
\begin{itemize}
  \item More than 90\% of the diffractive part comes from
    $x_\mathbb{P}>0.01$, at every rapidity. The $x_\mathbb{P}$ integral is
    dominated by its upper end, where the dipole is frozen.
  \item The form factor reduces the diffractive part by about one half, and
    more at negative rapidity. The ratio
    $\sigma(y=-2)/\sigma(y=-1)$ goes from $1.03$ to $0.91$. The maximum moves
    towards less negative $y$, as in the paper.
  \item The exclusive part changes by at most 9\% (at $y=-2$).
\end{itemize}

\section{What is not checked}

\begin{itemize}
  \item The test is only at $p_{D^0\perp}=2$~GeV. There are no
    $x_\mathbb{P}$-differential runs at 6--7 or 10--11~GeV, where the
    difference with the paper is largest. One run of \texttt{run\_xpom.sh}
    with \texttt{PT\_VALS="6.5 10.5"} and $y$ from $-2$ to $0$ would test it.
  \item Eq.~\eqref{eq:FA} is a simple form factor applied to the final
    integrand. It is not the factor of the paper, which is inside the nuclear
    PDF together with the shadowing.
  \item The form factor does not explain the difference in size: at high
    $p_T$ our result is already 5--10 times below the paper, and the form
    factor makes it smaller.
  \item Whether a coherence factor belongs in our formula is a physics
    decision. It is not in the code.
\end{itemize}

\end{document}
